\chapter{Open Challenges and Future Directions}
\label{sec:open-challenges}


% Springer Nature: each chapter should have a 10--15 line online abstract.
% This chapter abstract has not yet been written in the uploaded manuscript.
% [October 2026 round] Abstract written (blue).
\abstract{\added{This chapter collects the research problems that remain open once the classical, learning-based, hybrid, and multi-agent fusion methods of the preceding chapters are applied to autonomous driving and to the case studies beyond it. Eight challenges are examined in turn: robustness under adverse and long-tail conditions; calibration drift, sensor failure, and degraded modalities; uncertainty estimation, conflict handling, and explainability; real-time embedded deployment; verification, validation, and safety assurance; sim-to-real transfer and continual adaptation; privacy, standardisation, and interoperability; and foundation models and end-to-end fusion. For each challenge the chapter states the problem precisely, identifies which tools of the book already address part of it (innovation and consistency tests, covariance intersection, evidence discounting, zonotopic enclosures, MCDM-based reliability weighting, learned uncertainty, hybrid monitors), explains what is still missing, and formulates concrete open research questions. A summary table relates challenges, tools, and maturity, and the chapter closes with a consolidated research agenda and with discussion and computational problems on calibration error and on the effect of timestamp offsets.}}
\suppressfloats[t]

\begin{addedblock}
The preceding chapters have argued that sensor fusion is best understood as a single discipline with two traditions: the classical tradition of Chapter~\ref{sec:classicalSF}, which represents uncertainty explicitly and offers consistency checks, bounded-error enclosures, and conflict measures; and the learning tradition of Chapter~\ref{sec:learning}, which supplies recognition performance that no hand-designed model can match. Chapters~\ref{sec:application}--\ref{sec:applications} showed how the two interlock in autonomous driving, in multi-agent settings, and in two safety-critical domains beyond driving. This chapter asks what is still missing. For each of the eight challenges that follow we state the problem as precisely as the current state of knowledge allows, list the tools of the book that already address part of it, identify the gap, and formulate two or three open research questions. Table~\ref{tab:challenges-tools} summarises the mapping between challenges, tools, and maturity; Sect.~\ref{subsec:research-agenda} consolidates the questions into a research agenda. The problems at the end of the chapter are mostly design and discussion exercises, complemented by two computational ones (calibration error of a set of detection scores and the effect of a timestamp offset on an innovation test) whose scripts are in the companion repository.

% [October 2026 round] Summary table requested by the editor; the \color line is only needed while \revisionmarkuptrue.
\begin{table}[!t]
\ifrevisionmarkup\color{RevisionBlue}\fi
\caption{Open challenges, the tools of this book that already address part of each, and the maturity of the available solutions (\cmark\ mature and deployable, \pmark\ partial or active research, \xmark\ largely open).}
\label{tab:challenges-tools}
\footnotesize
\setlength{\tabcolsep}{3pt}
\begin{tabular}{>{\raggedright\arraybackslash}p{0.21\textwidth} >{\raggedright\arraybackslash}p{0.40\textwidth} c >{\raggedright\arraybackslash}p{0.24\textwidth}}
\toprule
Challenge & Relevant tools in the book & Maturity & Main gap \\
\midrule
Adverse and long-tail robustness (Sect.~\ref{subsec:oc-robustness}) & Innovation/NIS test \eqref{eq:nis}; evidence conflict and discounting (Chapter~\ref{sec:classicalSF}); MCDM reliability weights; adverse-weather benchmarks (Chapter~\ref{sec:application}) & \pmark & Inconsistency is detected only after the fact; no predictive model of modality-specific degradation \\[2pt]
Calibration drift and sensor failure (Sect.~\ref{subsec:oc-calibration}) & NIS monitoring; augmented-state (online) calibration; evidence discounting; zonotopes with inflated bounds & \pmark & Identifiability of drift versus model mismatch; guaranteed detection delay \\[2pt]
Uncertainty, conflict, explainability (Sect.~\ref{subsec:oc-uncertainty}) & Learned covariance as $\mathbf{R}$ (Fig.~\ref{fig:hybrid-pipeline}); conflict mass $K$ \eqref{eq:dempster}; pignistic transform; fuzzy and explainable fusion (Chapter~\ref{sec:learning}) & \pmark & Calibration under distribution shift; conversion between uncertainty representations \\[2pt]
Real-time embedded deployment (Sect.~\ref{subsec:oc-embedded}) & FPGA estimation cores (Chapter~\ref{sec:application}); efficient DNN processing; quantisation and distillation & \cmark/\pmark & Effect of compression on calibration; anytime fusion under compute contention \\[2pt]
Verification, validation, safety (Sect.~\ref{subsec:oc-vv}) & Explicit monitor signals; set-membership guarantees; HIL and infrastructure testbeds (Chapter~\ref{sec:application}); scenario-based testing & \xmark & Coverage metrics for fusion failures; composition of assurance; validation of learned uncertainty \\[2pt]
Sim-to-real and continual adaptation (Sect.~\ref{subsec:oc-sim2real}) & Per-modality measurement models (Chapters~\ref{sec:Principles}, \ref{sec:classicalSF}); NIS as a transfer metric; domain randomisation and adaptation & \pmark & Per-modality transfer metrics; re-validation after adaptation \\[2pt]
Privacy, standardisation, interoperability (Sect.~\ref{subsec:oc-privacy}) & Covariance intersection \eqref{eq:ci-est}; consensus \eqref{eq:consensus-filter}; discounting (Chapter~\ref{sec:MAS}); federated learning & \xmark & Uncertainty and provenance in message standards; trust models \\[2pt]
Foundation models and end-to-end fusion (Sect.~\ref{subsec:oc-foundation}) & Discounting of an additional source; hybrid monitors (Chapter~\ref{sec:learning}) & \xmark & Location of the safety layer; calibrated uncertainty from latent representations \\
\bottomrule
\end{tabular}
\end{table}
\end{addedblock}

\section{Robustness Under Adverse and Long-Tail Conditions}
\label{subsec:oc-robustness}

Sensor-fusion systems must operate under rare, degraded, and difficult-to-model conditions. In autonomous driving, these include adverse weather, unusual lighting, rare road users, roadworks, emergency vehicles, sensor blockage, and unusual traffic behaviour. In other domains, such as maritime monitoring or critical-infrastructure protection, long-tail events may include spoofed reports, missing data, rare vessel behaviours, or infrequent physical disturbances. Public datasets only partially cover such cases, while synthetic augmentation may not fully reproduce real sensor degradation. Future work should develop more realistic sensor-corruption models, richer adverse-condition datasets~\cite{bijelic2020_stf,dong2023_robobev}, and evaluation protocols that reward robust behaviour rather than only average-case performance.

\begin{addedblock}
Stated precisely, the challenge is one of distribution shift with modality-specific structure. A fusion system is designed, tuned, and validated under a nominal distribution of sensor returns, but it must operate under conditions whose probability per hour of driving is far below $10^{-6}$ and whose effect on each modality is different: glare, low light, and lens contamination affect cameras; spray, fog backscatter, and wet-surface absorption affect LiDAR; multipath, clutter, and interference affect radar. The consequence for fusion is that the sensor weights, covariances, and reliability factors that were appropriate under nominal conditions become wrong in a correlated way, and the fused estimate can be confidently wrong rather than merely uncertain. Benchmarks that corrupt inputs systematically, in the style of common-corruption suites for image classifiers \cite{hendrycks2019_corruptions} and their extensions to BEV perception \cite{dong2023_robobev}, confirm that the degradation of multimodal networks is non-uniform across modalities and corruption types, and that models trained in clear weather can lose a large fraction of their accuracy on the exact subsets for which fusion was introduced \cite{bijelic2020_stf,feng2021_multimodal_survey}.

The book already supplies several instruments. The normalised innovation squared \eqref{eq:nis} and the track-level gates of Chapter~\ref{sec:classicalSF} detect that a sensor has become inconsistent with the filter model; evidence theory quantifies the conflict between sources through the mass $K$ in Dempster's rule and lets the designer discount an unreliable source rather than discard it; the MCDM-based weighting of Chapter~\ref{sec:classicalSF} makes reliability weights context dependent, so that a rain or night context can be mapped to a different weight vector; and the set-membership estimators provide enclosures that remain valid whenever the bounded-error assumption holds, regardless of the shape of the noise. Chapter~\ref{sec:application} reviews the datasets that contain adverse-weather subsets. What is missing is predictive rather than reactive robustness: all of the above detect inconsistency after it has occurred, whereas a robust fusion system should anticipate it from context (weather, time of day, sensor self-diagnostics) and adjust before the first wrong estimate reaches the planner. There is, moreover, no accepted metric that rewards graceful degradation: average precision on a corrupted test set measures accuracy but not whether the system knew that its accuracy had dropped.

\begin{list4a}
    \item \textbf{Open question 1.} Can modality-specific degradation be modelled physically well enough that a context variable (precipitation rate, illuminance, interference level) predicts the covariance inflation or discount factor of each sensor, replacing the reactive adaptation of innovation-based schemes?
    \item \textbf{Open question 2.} Which evaluation protocol rewards robust behaviour, for example a joint score of accuracy and calibration under each corruption, or a penalty for confident errors, and can it be adopted by the public benchmarks of Chapter~\ref{sec:application}?
    \item \textbf{Open question 3.} How should long-tail events that are rare in every dataset, such as emergency vehicles or spoofed cooperative messages, be weighted during training and validation without biasing nominal performance?
\end{list4a}
\end{addedblock}

\section{Calibration Drift, Sensor Failure, and Degraded Modalities}
\label{subsec:oc-calibration}

Fusion performance depends strongly on calibration, synchronisation, and sensor health. Mechanical vibration, temperature change, sensor replacement, clock drift, communication delays, or physical obstruction can introduce spatial or temporal misalignment. Sensors may also be blocked, degraded, spoofed, saturated, or damaged. Future systems require online calibration monitoring, automatic fault detection, sensor-health estimation, and fusion architectures that degrade gracefully when one or more modalities become unreliable.

\begin{addedblock}
In the notation of Chapter~\ref{sec:classicalSF}, calibration enters the measurement model as a set of nuisance parameters: the extrinsic transform $\mathbf{T}_i$ between sensor $i$ and the vehicle frame and the time offset $\delta_i$ between its clock and the reference clock appear inside $h_i(\mathbf{x};\mathbf{T}_i,\delta_i)$. Drift is a slowly varying error in these parameters, and its effect on the fused estimate is a bias that is not represented in the covariance. The effect of a time offset is particularly insidious because it masquerades as a velocity-dependent position bias of magnitude $\lVert\mathbf{v}\rVert\,\delta_i$: a $20$~ms offset is invisible for a parked target and equals a $0.3$~m bias for a target closing at $15$~m/s (Problem~\ref{prob:ch08:3}). The practical state of the art estimates $\mathbf{T}_i$ and $\delta_i$ jointly, either offline with dedicated targets \cite{furgale2013_kalibr} or online by maximising the consistency of co-observed structure \cite{levinson2013_online_calibration}; the latter is the natural route to drift monitoring, since the same objective can be evaluated continuously.

The tools of the book cover the detection side reasonably well. The NIS test \eqref{eq:nis} and the windowed consistency statistics of Chapter~\ref{sec:classicalSF} flag a drifting sensor once its bias exceeds a fraction of the innovation standard deviation; augmented-state estimation, in which $\mathbf{T}_i$ and $\delta_i$ are appended to the state vector, converts drift into an estimable quantity; evidence discounting and MCDM weights provide a principled way to reduce the influence of a suspect sensor; and zonotopic estimation with inflated bounds (Chapter~\ref{sec:classicalSF}) keeps a guaranteed enclosure while the fault is diagnosed. The hybrid pipeline of Chapter~\ref{sec:learning} (Fig.~\ref{fig:hybrid-pipeline}) shows where these monitors sit relative to the learned detector. What is missing is identifiability and timing. A persistent innovation bias can be caused by calibration drift, by a wrong process model, by a degraded sensor, or by a data-association error, and the NIS alone cannot tell which; nor does the literature offer guaranteed bounds on the detection delay of such monitors as a function of the drift rate. Finally, ``graceful degradation'' is usually described qualitatively, whereas a safety argument needs it as a quantitative design objective: which fused-accuracy level is guaranteed when sensor $i$ is lost, and for how long.

\begin{list4a}
    \item \textbf{Open question 1.} Which combination of residual signatures (temporal correlation of innovations, dependence on ego velocity, cross-sensor consistency) makes drift, degradation, process-model mismatch, and association errors mutually identifiable, and with what minimum detection delay?
    \item \textbf{Open question 2.} How should the fusion architecture be specified so that the loss or discounting of any single modality yields a provable, pre-computed degradation of accuracy and integrity, in the sense of a protection level?
    \item \textbf{Open question 3.} Can learned perception be made self-diagnosing, for example by predicting its own calibration state or by exposing features whose statistics drift measurably with sensor health?
\end{list4a}
\end{addedblock}

\section{Uncertainty Estimation, Conflict Handling, and Explainability}
\label{subsec:oc-uncertainty}

Many learning-based fusion systems produce accurate predictions but poorly calibrated confidence estimates. For safety-critical applications, the system must know when it does not know. Future research should improve uncertainty estimation, out-of-distribution detection, explainable sensor weighting, and conflict detection between modalities. Hybrid approaches that combine neural perception with probabilistic estimation, evidence theory, fuzzy logic, or interpretable reliability models are promising because they can expose why a sensor was trusted, down-weighted, or treated as conflicting.

\begin{addedblock}
The precise requirement is calibration: a confidence score $p$ attached to a detection or a class label is calibrated if, among all outputs with confidence $p$, a fraction $p$ is correct. The standard summary is the expected calibration error, $\mathrm{ECE}=\sum_{b=1}^{M}\frac{n_b}{N}\,\lvert\mathrm{acc}_b-\mathrm{conf}_b\rvert$, computed over $M$ confidence bins (Problem~\ref{prob:ch08:2}). Modern deep networks are systematically over-confident, and simple post-hoc remedies such as temperature scaling recover calibration on in-distribution data \cite{guo2017_calibration}. Uncertainty itself has two sources that must be handled differently: aleatoric uncertainty, intrinsic to the measurement, which can be learned as a heteroscedastic output and used directly as a covariance; and epistemic uncertainty, due to limited training data, which requires Bayesian approximations such as Monte Carlo dropout \cite{gal2016_mcdropout} or deep ensembles \cite{lakshminarayanan2017_ensembles,kendall2017_uncertainties}. For LiDAR-based 3D detection, capturing both has been shown to improve the usefulness of detections downstream \cite{feng2018_lidar_uncertainty}. Out-of-distribution detection, starting from the maximum-softmax baseline \cite{hendrycks2017_ood}, addresses the complementary question of whether the input is one on which any confidence can be trusted. The central difficulty for fusion is that calibration does not survive distribution shift: confidence estimates that are calibrated on the training distribution degrade, often in the over-confident direction, precisely under the corruptions of Sect.~\ref{subsec:oc-robustness} \cite{ovadia2019_shift}.

The book contributes the receiving side of this interface. In the hybrid pipeline of Chapter~\ref{sec:learning}, a learned covariance enters the tracker as the measurement-noise matrix $\mathbf{R}$, so that a poorly calibrated network is immediately visible as an inconsistent NIS sequence, and the filter residuals provide supervision-free feedback on calibration. Evidence theory (Chapter~\ref{sec:classicalSF}) offers a representation in which ignorance is explicit: a detector that does not know can place mass on the frame of discernment rather than on a class, the conflict mass $K$ of \eqref{eq:dempster} quantifies disagreement between modalities, Smets' transferable belief model \eqref{eq:smets} keeps conflict visible instead of renormalising it away, and the pignistic transform \eqref{eq:pignistic-transform} converts beliefs back into decisions. The fuzzy and explainable fusion schemes of Chapter~\ref{sec:learning} and the MCDM weights of Chapter~\ref{sec:classicalSF} expose why a sensor was trusted. What is missing is a common currency. Covariances, credal sets, basic belief assignments, zonotopes, and softmax scores are different representations of uncertainty, and the conversions between them (for instance from a learned covariance to a discount factor, or from a conflict mass to a covariance inflation) are chosen ad hoc. Explainability is likewise asserted rather than measured: there is no agreed test of whether the explanation offered by a reliability model (``radar was down-weighted because of conflict with LiDAR'') is faithful to the computation that produced the estimate.

\begin{list4a}
    \item \textbf{Open question 1.} Can calibration be maintained under distribution shift by using the innovation statistics of the classical layer as an online recalibration signal, closing the loop sketched in Fig.~\ref{fig:hybrid-pipeline} with provable guarantees?
    \item \textbf{Open question 2.} What is a principled mapping between the uncertainty representations used in this book (covariance, bounded set, basic belief assignment, fuzzy membership) such that a single reliability signal can be propagated across the layers of a hybrid pipeline without loss of meaning?
    \item \textbf{Open question 3.} How should the faithfulness of a fusion explanation be evaluated, and can explanation quality be made a requirement that a safety case can reference?
\end{list4a}
\end{addedblock}

\section{Real-Time Embedded Deployment}
\label{subsec:oc-embedded}

Modern fusion architectures can be computationally expensive, especially when they use high-resolution images, dense point clouds, radar tensors, distributed acoustic measurements, temporal histories, or transformer-based attention. Real-time deployment requires careful optimisation of latency, memory, power consumption, communication bandwidth, and hardware utilisation, on automotive-grade platforms such as NVIDIA DRIVE Orin, Qualcomm Snapdragon Ride, TI TDA4-class SoCs, and FPGA-based accelerators%
~\cite{lin2018_architectural,liu2019_edge_computing,sze2017_efficient}.
Future research should consider efficient model design, sparse computation, quantisation, hardware-aware training, and co-design between algorithms and embedded platforms.

\begin{addedblock}
Latency deserves to be treated as a fusion error rather than as an engineering nuisance. If the fused estimate of an object is delivered $\tau$ seconds after the measurements were taken, the planner acts on a position that is wrong by $\lVert\mathbf{v}\rVert\,\tau$ for a target with relative velocity $\mathbf{v}$: this is the same mechanism as a timestamp offset (Sect.~\ref{subsec:oc-calibration}), except that it is deterministic and known, and can therefore be compensated by prediction provided that the motion model is adequate. The compute budget of a fusion stack is thus coupled to its accuracy requirement, and the usual end-to-end budget of roughly $100$~ms for perception and fusion at $10$~Hz is not arbitrary. Within that budget, dense BEV fusion networks and transformer-based attention over multiple cameras and LiDAR sweeps are the dominant cost; the classical layer (filters, gating, covariance intersection, evidence combination) is cheap by comparison, which is one more argument for keeping it explicit. Established efficiency techniques are pruning and compression \cite{han2016_deep_compression}, integer quantisation with quantisation-aware training \cite{jacob2018_quantization}, and knowledge distillation into smaller students \cite{hinton2015_distillation}; the architectural implications of running such networks on automotive SoCs and edge devices are surveyed in \cite{lin2018_architectural,liu2019_edge_computing,sze2017_efficient}.

The book covers the classical side of deployment in Chapter~\ref{sec:application}, where FPGA implementations of Kalman and LQG cores illustrate how estimation can be accelerated and validated with hardware-in-the-loop infrastructure, and in Chapter~\ref{sec:MAS}, where communication bandwidth and latency constrain which information-exchange level is feasible. What is missing is an understanding of how compression interacts with the uncertainty interface of Sect.~\ref{subsec:oc-uncertainty}: quantisation and pruning are evaluated on accuracy, almost never on calibration, so a compressed detector may deliver a covariance that is no longer consistent with its errors. Equally open is scheduling: on a shared SoC the fusion task competes with perception, prediction, and planning, and a fusion algorithm whose quality degrades gracefully when its time slice is cut (an anytime algorithm) would be preferable to one that simply misses its deadline.

\begin{list4a}
    \item \textbf{Open question 1.} How do quantisation, pruning, and distillation affect the calibration of learned uncertainty, and can calibration be preserved by construction during hardware-aware training?
    \item \textbf{Open question 2.} Can fusion be formulated as an anytime computation, for example by ordering sensors or features by information gain so that an early deadline yields a valid, conservatively bounded estimate?
    \item \textbf{Open question 3.} What is the right division of a fusion stack between a learned part that benefits from accelerators and a classical part that benefits from deterministic timing, and how is timing determinism itself argued in a safety case?
\end{list4a}
\end{addedblock}

\section{Verification, Validation, and Safety Assurance}
\label{subsec:oc-vv}

Sensor-fusion systems must be validated under a very large scenario space, including rare events, sensor failures, degraded sensing, and conflicting inputs. Simulation, scenario-based testing, closed-course testing, hardware-in-the-loop systems, real-world trials, and formal safety arguments are all required. A key open problem is how to validate learning-based fusion systems whose behaviour depends on high-dimensional data and may not be easily interpretable. Validation should therefore assess not only task accuracy, but also uncertainty calibration, fault tolerance, graceful degradation, and the effect of fusion errors on downstream decisions. Emerging automotive standards give this challenge a concrete regulatory shape: ISO~26262 addresses functional safety of electronic systems, ISO~21448 (SOTIF) addresses hazards arising from functional insufficiencies of the intended behaviour --- including perception and fusion errors without any hardware fault --- and ISO/PAS~8800 targets safety of AI-based functions in road vehicles. A fusion architecture whose uncertainty, conflict, and degradation signals are explicit, as advocated throughout this survey, is substantially easier to argue about under these frameworks than an end-to-end black box. % [July 2026 pass] A V&V subsection in a 2026 automotive survey should acknowledge the applicable safety standards; added without new bib entries (standards are cited by designation, as is common in IEEE Access).

\begin{addedblock}
The size of the problem can be stated statistically. To demonstrate with $95\%$ confidence that a fleet's fatality rate is below the human benchmark by driving alone would require hundreds of millions of miles, and comparing two system versions would require billions \cite{kalra2016_miles}; statistical validation by mileage is therefore not a viable strategy on its own. The alternatives are decomposition and argument: scenario-based testing concentrates the test budget on a structured catalogue of operational situations \cite{riedmaier2020_scenario}, formal models such as responsibility-sensitive safety constrain the planner so that perception is the only source of residual risk \cite{shalevshwartz2017_rss}, and the engineering challenges of testing machine-learning components inside such a decomposition are discussed in \cite{koopman2016_testing} and, from the machine-learning side, in \cite{amodei2016_concrete}. Fusion occupies a specific place in this decomposition: it is the component that converts raw perception errors into the state estimates and uncertainties on which the planner's safety argument depends, so its validation must address not only accuracy but the correctness of its uncertainty, exactly the SOTIF notion of functional insufficiency.

The tools of the book help in two ways. First, every classical mechanism exposes a signal that can become a monitorable requirement: the NIS sequence must remain within its $\chi^2$ bounds for a specified fraction of time, the conflict mass must stay below a threshold, the zonotopic enclosure must contain the true state whenever the noise bounds hold. Requirements of this form can be tested in simulation, in hardware-in-the-loop set-ups, and on infrastructure testbeds such as the one described in Chapter~\ref{sec:application}, and their violation rate can be reported as evidence. Second, set-membership estimation provides a worst-case guarantee under bounded noise, which is the type of statement a safety case can use directly, whereas probabilistic estimates require an additional argument about the validity of the noise model. What is missing is coverage and composition. There is no accepted coverage metric for fusion-specific failure modes (association errors, correlated sensor degradation, delayed or out-of-order messages), so it is not possible to say when scenario-based testing of a fusion module is complete; and there is no established method for composing component-level evidence (a calibrated detector, a consistent filter, a conflict monitor) into a system-level claim about the planner's residual risk. Validating learned uncertainty is itself open: the calibration metrics of Sect.~\ref{subsec:oc-uncertainty} are computed on datasets, not on the operational design domain.

\begin{list4a}
    \item \textbf{Open question 1.} What coverage criteria are appropriate for the fusion layer, for example a catalogue of sensor-degradation and association-failure modes with required test outcomes, and how do they map onto ISO~21448 triggering conditions?
    \item \textbf{Open question 2.} How can the monitor signals of the classical layer (NIS, conflict, enclosure containment) be turned into run-time assurance arguments acceptable under ISO/PAS~8800, including bounds on their own false-alarm and missed-detection rates?
    \item \textbf{Open question 3.} Can the residual risk contributed by fusion errors be propagated analytically to the planner, so that the required fusion integrity follows from the planner's safety envelope rather than being set by convention?
\end{list4a}
\end{addedblock}

\section{Sim-to-Real Transfer and Continual Adaptation}
\label{subsec:oc-sim2real}

% % Addressed in the previous round: subsection rewritten (blue block).

Learning-based fusion models are increasingly trained or augmented in simulation, because rare and dangerous scenarios are cheap to generate synthetically. However, simulated sensor models rarely reproduce the detailed failure characteristics of real cameras, LiDAR, and radar --- noise textures, multipath, blooming, spray, or weather-dependent attenuation --- so models can inherit a sim-to-real gap precisely in the degraded conditions for which simulation was used. Closing this gap requires physically faithful sensor simulation, domain-adaptation techniques evaluated per modality, and validation protocols that measure transfer explicitly. A related challenge is continual and online learning: fusion systems that adapt during deployment (recalibrating reliability weights, updating maps, or fine-tuning detectors) must do so without catastrophic forgetting and without invalidating prior safety arguments.

\begin{addedblock}
Two families of techniques address the gap from the learning side. Domain randomisation deliberately varies the rendering parameters of the simulator so that the real world appears to the network as one more variation \cite{tobin2017_domain_randomization}; domain adaptation instead aligns the feature or pixel distributions of the simulated and real domains, for instance with cycle-consistent image translation \cite{hoffman2018_cycada}. Both have been developed primarily for cameras, and the simulators used in driving research \cite{dosovitskiy_carla_2017,shah_airsim_2017} are also camera-centric; data-driven LiDAR simulation that composes real scans into new scenes narrows the gap for one further modality \cite{manivasagam2020_lidarsim}, while faithful radar simulation remains the least developed. The consequence for fusion is that the sim-to-real gap differs by modality, so a fusion network trained in simulation may learn modality weights that are correct in simulation and wrong in reality. On the continual-learning side, regularisation methods such as elastic weight consolidation protect parameters that were important for earlier tasks \cite{kirkpatrick2017_ewc}, and the broader landscape of continual learning is reviewed in \cite{parisi2019_continual}; what none of these methods addresses is the safety-assurance question of whether a model that has changed is still the model that was validated.

The book offers a vocabulary for measuring the gap rather than closing it. Chapters~\ref{sec:Principles} and~\ref{sec:classicalSF} describe the per-modality measurement models against which a simulator's output can be compared statistically, and the innovation statistics of Chapter~\ref{sec:classicalSF} provide a transfer metric that needs no labels: a detector trained in simulation and deployed on real data will exhibit an inconsistent NIS sequence if its learned covariances are mis-scaled, and the ratio of real to simulated NIS is a direct measure of how much of the uncertainty model transferred (Problem~\ref{prob:ch08:4}). The adaptive reliability weights of Chapters~\ref{sec:classicalSF} and~\ref{sec:learning} are the part of a fusion system that can be updated online most safely, since they are low dimensional and interpretable. What is missing is a per-modality transfer protocol with accept/reject criteria, and an understanding of which parameters of a fusion system may be adapted online without re-validation: a continual-learning procedure that preserves accuracy but silently degrades calibration would pass every benchmark and fail every safety argument.

\begin{list4a}
    \item \textbf{Open question 1.} Which label-free statistics (innovation consistency, conflict rate, calibration on pseudo-labels) are sufficient to certify that a simulation-trained fusion component transfers, modality by modality, and what thresholds should trigger re-training?
    \item \textbf{Open question 2.} Can a fusion architecture be partitioned into a frozen, validated core and an adaptable periphery (reliability weights, discount factors, covariance scalings) so that continual adaptation is confined to parameters whose effect on safety is bounded?
    \item \textbf{Open question 3.} How should physically faithful radar and LiDAR degradation models be validated against real adverse-weather recordings so that they can be trusted to generate the long-tail cases of Sect.~\ref{subsec:oc-robustness}?
\end{list4a}
\end{addedblock}

\section{Privacy, Standardisation, and Interoperability}
\label{subsec:oc-privacy}
% % Addressed in the previous round: subsection rewritten (blue block).
Multi-agent and infrastructure-assisted fusion make data governance a first-order design constraint. Cooperative perception exchanges observations of identifiable people and vehicles across organisational boundaries, raising privacy, retention, and consent questions that pure onboard fusion avoids; techniques such as feature-level sharing, anonymisation, and federated learning trade information content against privacy exposure. Interoperability is equally unresolved: V2X message formats, coordinate conventions, time bases, uncertainty representations, and trust models must be standardised before fused estimates from heterogeneous vendors can be combined safely. Without agreed ways to exchange covariances, confidences, and provenance, receiving agents cannot apply the consistency mechanisms (covariance intersection, evidence discounting) that this survey has argued are essential.

\begin{addedblock}
The three issues are coupled through the information-exchange level of Chapter~\ref{sec:MAS}. Raw-data sharing is the most informative and the least private; track-level sharing, in which an agent transmits object states with covariances, is compact and is what existing standards foresee, with ETSI's collective perception message (ETSI~TR~103~562) and the SAE~J2735 message set defining the fields; feature-level sharing, in which intermediate neural representations are exchanged and fused by the receiver, has been shown to recover most of the benefit of raw sharing at a fraction of the bandwidth \cite{wang2020_v2vnet} and now has open benchmarks \cite{xu2022_opv2v}, but its privacy properties are poorly understood because learned features can leak identifiable content. Federated learning \cite{mcmahan2017_fedavg} removes the need to share data at all when the goal is to improve a shared model rather than a shared estimate, at the cost of open problems in heterogeneity, robustness, and verification \cite{kairouz2021_federated}; differential privacy \cite{dwork2006_dp,dwork2014_algorithmic_dp} offers the only formal privacy guarantee available, but its noise addition is in direct tension with the accuracy that fusion is meant to deliver.

The book's contribution is to make explicit what a receiving agent needs. Covariance intersection \eqref{eq:ci-est}--\eqref{eq:ci-cov} requires a covariance; evidence discounting requires a reliability factor or a basis for estimating one; consensus \eqref{eq:consensus-filter} requires a common state definition, frame, and time base; and all of them require provenance, so that the same measurement is not fused twice through different paths (the rumour-propagation problem of Chapter~\ref{sec:MAS}). None of these items is fully specified by current message standards, which carry position and a scalar confidence at best; a receiver therefore has to guess a covariance, and the consistency guarantees of Chapter~\ref{sec:classicalSF} are lost (Problem~\ref{prob:ch08:5}). What is missing is, first, a standardised uncertainty and provenance payload with a defined semantics; second, a trust model that lets a receiver discount a sender whose declared covariance has proven inconsistent, which the innovation test of the receiver can detect; and third, a quantitative understanding of how differential-privacy noise or feature anonymisation propagates into fused accuracy and integrity.

\begin{list4a}
    \item \textbf{Open question 1.} What minimal uncertainty and provenance payload (covariance or bounded set, time base, frame, sensor lineage, declared reliability) makes consistent fusion possible across vendors, and how can its correctness be audited at run time?
    \item \textbf{Open question 2.} How should a receiving agent learn trust in a sender from the consistency of its messages, and how is that trust combined with evidence discounting without creating an attack surface for spoofing?
    \item \textbf{Open question 3.} What is the achievable trade-off between a formal privacy guarantee on shared perception and the accuracy and integrity of the resulting cooperative estimate?
\end{list4a}
\end{addedblock}

\section{Foundation Models and End-to-End Fusion}
\label{subsec:oc-foundation}
%
Recent progress in large-scale representation learning and foundation models may influence future sensor-fusion systems. Large multimodal models could improve generalisation, semantic reasoning, and data efficiency, while end-to-end architectures may learn integrated perception, prediction, planning, or monitoring representations. However, these approaches raise important questions about interpretability, safety, uncertainty, data bias, compute requirements, and failure analysis. Future fusion research must determine how to combine the flexibility of large learned models with the reliability, structure, and auditability required for safety-critical operation.

Two concrete developments illustrate the trend. First, planning-oriented end-to-end architectures such as UniAD train perception, prediction, and planning jointly in a single network, so that fusion is no longer an explicit module but an emergent property of intermediate representations \cite{hu2023_uniad}; recent surveys of end-to-end driving systematise the rapidly growing design space and its open problems, including causal confusion, covariate shift, and the difficulty of interpreting intermediate features \cite{chen2024_e2e_survey}. Second, large multimodal language models are being explored as high-level reasoning layers on top of fused perception, for instance to interpret rare scenes, explain decisions in natural language, or serve as a semantic prior for long-tail objects \cite{cui2024_mllm_driving}. For the fusion perspective taken in this survey, the key open questions are where the classical safety layer should sit when the fusion function is implicit --- for example, whether innovation-style consistency checks can be applied to latent representations --- how calibrated uncertainty can be extracted from end-to-end models, and whether foundation-model outputs can be treated as one more (highly capable but occasionally unreliable) information source to be discounted and audited with the conflict-handling machinery of Chapter~\ref{sec:classicalSF}.

\begin{addedblock}
The precise difficulty is that the unified view developed in this book relies on an explicit interface between perception and estimation: measurements with covariances enter a filter, beliefs with masses enter a combination rule, sets with bounds enter an intersection. An end-to-end network removes this interface, and the properties that made the classical layer useful (consistency checks, explicit conflict, enclosures) have no obvious counterpart in a latent representation. The risks are documented: imitation-trained driving policies can suffer from causal confusion, learning to rely on spurious correlates in the input \cite{dehaan2019_causal_confusion}; and foundation models as a class bring homogenisation of failure modes and difficulties of evaluation that have been catalogued at length \cite{bommasani2021_foundation}. At the same time, the planning literature \cite{schwarting2018_planning} makes clear that planners already consume uncertainty in structured form, so the question is not whether a structured interface is needed but where to place it.

The book suggests two placements that are compatible with the tools of Chapters~\ref{sec:classicalSF} and~\ref{sec:learning}. The first treats the foundation model or end-to-end network as one more information source: its outputs (object hypotheses, semantic labels, scene descriptions) are mapped to basic belief assignments with a discount factor estimated from held-out performance, combined with the explicit sensors through the rules of Chapter~\ref{sec:classicalSF}, and audited through the conflict mass (Problem~\ref{prob:ch08:6}). The second keeps a classical estimator in parallel with the end-to-end system as an independent reference, so that the innovation between the two becomes a consistency test on the learned pipeline; this is the architecture of a run-time monitor, and it requires only that the end-to-end system expose some estimable quantities (ego pose, nearest-object states). What is missing is a theory of consistency checking for latent representations, a way to extract calibrated uncertainty from models that were never trained to produce it, and benchmarks that evaluate the combination of a foundation model with a classical safety layer rather than either alone.

\begin{list4a}
    \item \textbf{Open question 1.} Can innovation-style consistency tests be defined on latent representations, for example by learning a decoder to a small set of physical quantities that a classical estimator can predict, and what statistical guarantees would such tests carry?
    \item \textbf{Open question 2.} How should the discount factor of a foundation-model source be estimated and updated online, given that its failure modes are correlated across scenes and poorly characterised?
    \item \textbf{Open question 3.} Which architecture, a foundation model audited by a classical layer or a classical pipeline enriched by foundation-model priors, offers the better trade-off between long-tail generalisation and assurability, and under which operational design domains?
\end{list4a}

\subsection{A Consolidated Research Agenda}
\label{subsec:research-agenda}

The questions above recur across sections, and it is useful to close with the themes that connect them. Each item below names a research direction, the chapters whose tools it builds on, and the challenge it would resolve.
\begin{list4a}
    \item \textbf{Predictive reliability models.} Replace reactive innovation-based adaptation by context-driven prediction of per-modality covariance inflation and discount factors (Chapters~\ref{sec:classicalSF} and~\ref{sec:learning}; Sects.~\ref{subsec:oc-robustness} and~\ref{subsec:oc-calibration}).
    \item \textbf{A common currency for uncertainty.} Principled mappings between covariances, bounded sets, belief masses, and learned confidences, so that a single reliability signal propagates through a hybrid pipeline (Chapter~\ref{sec:classicalSF}; Sect.~\ref{subsec:oc-uncertainty}).
    \item \textbf{Calibration as a first-class requirement.} Benchmarks, compression methods, and continual-learning procedures that are evaluated on calibration under shift and not only on accuracy (Chapters~\ref{sec:learning} and~\ref{sec:application}; Sects.~\ref{subsec:oc-uncertainty}, \ref{subsec:oc-embedded}, and~\ref{subsec:oc-sim2real}).
    \item \textbf{Monitor signals as assurance evidence.} Turning NIS bounds, conflict thresholds, and enclosure containment into run-time requirements with characterised false-alarm and missed-detection rates, mapped onto ISO~21448 and ISO/PAS~8800 (Chapters~\ref{sec:classicalSF} and~\ref{sec:application}; Sect.~\ref{subsec:oc-vv}).
    \item \textbf{Frozen core, adaptable periphery.} Architectures in which online adaptation is confined to low-dimensional, interpretable parameters whose safety effect is bounded (Chapter~\ref{sec:learning}; Sect.~\ref{subsec:oc-sim2real}).
    \item \textbf{Uncertainty-carrying cooperative standards.} Message payloads with covariance, provenance, time base, and declared reliability, together with trust models that receivers can audit (Chapter~\ref{sec:MAS}; Sect.~\ref{subsec:oc-privacy}).
    \item \textbf{Classical layers for implicit fusion.} Consistency tests and discounting schemes that apply to end-to-end and foundation-model outputs, evaluated on benchmarks that score the combination rather than the parts (Chapters~\ref{sec:classicalSF} and~\ref{sec:learning}; Sect.~\ref{subsec:oc-foundation}).
    \item \textbf{Transfer to other domains.} The maritime and submarine-cable case studies of Chapter~\ref{sec:applications} show that every item above has a counterpart outside driving; a research programme that validates the same tools in several domains will produce more robust answers than one confined to road vehicles.
\end{list4a}
\end{addedblock}

\begin{addedblock}
\problemsection

\begin{prob}\label{prob:ch08:1}
\textbf{Sensor-health monitor (design).} A vehicle fuses camera, radar, and LiDAR detections in a Kalman tracker as in Fig.~\ref{fig:hybrid-pipeline}. Design a sensor-health monitor that uses the per-sensor NIS \eqref{eq:nis}, an evidence-theoretic discount factor per sensor, and a graceful-degradation policy. Specify (a) the statistics that are computed and over which window, (b) the thresholds and how their false-alarm rate is set, (c) how the discount factor is derived from the statistics, and (d) how the policy distinguishes a blocked sensor (no detections) from a miscalibrated one (biased detections) and what the planner is told in each case.
\end{prob}

\begin{prob}\label{prob:ch08:2}
\textbf{Expected calibration error (computational).} A detector reports the following twelve confidences, and ground truth states whether each detection was correct ($1$) or not ($0$):
\begin{center}\scriptsize
\setlength{\tabcolsep}{3.5pt}
\begin{tabular}{l*{12}{c}}
\toprule
$p_i$ & $0.95$ & $0.92$ & $0.88$ & $0.85$ & $0.78$ & $0.72$ & $0.65$ & $0.58$ & $0.52$ & $0.45$ & $0.35$ & $0.22$ \\
$y_i$ & $1$ & $1$ & $0$ & $1$ & $1$ & $0$ & $1$ & $0$ & $0$ & $0$ & $1$ & $0$ \\
\bottomrule
\end{tabular}
\end{center}
(a) Compute the expected calibration error with $M=5$ equal-width bins on $[0,1]$. (b) Compute the Brier score $\frac{1}{N}\sum_i(p_i-y_i)^2$ and the overall accuracy and mean confidence. (c) Is the detector over- or under-confident? Temperature scaling replaces $p_i$ by $\sigma(z_i/T)$ with $z_i=\log\frac{p_i}{1-p_i}$; state whether $T>1$ or $T<1$ is needed and verify numerically by finding the $T$ that minimises the negative log-likelihood, using the companion script \repopath{ch08\_\allowbreak p2\_\allowbreak ece.py} in \repopath{companion\_\allowbreak repo/\allowbreak exercises/\allowbreak ch08/\allowbreak solutions/\allowbreak }. (d) Comment on the reliability of an ECE estimate from twelve samples.
\end{prob}

\begin{prob}\label{prob:ch08:3}
\textbf{Timestamp offset and the innovation test (computational).} A radar measures the range to a lead vehicle with standard deviation $\sigma=0.3$~m, and the measurement is fused in a filter whose innovation covariance is $S=\sigma^2$ (scalar measurement, negligible prediction uncertainty). The radar clock is offset by an unknown $\delta$ with respect to the fusion clock, so that each measurement is interpreted $\delta$ seconds late. The lead vehicle closes at relative speed $v=15$~m/s. (a) Show that the offset produces an innovation bias $b=v\delta$ and compute $b$ for $\delta=20$~ms and $\delta=50$~ms. (b) Show that the expected NIS \eqref{eq:nis} becomes $1+b^2/\sigma^2$ and evaluate it for both offsets. (c) The single-sample $95\%$ gate for $\chi^2_1$ is $3.84$: for which of the two offsets is a correct measurement rejected on average? (d) A windowed test averages the NIS over $N=20$ samples and compares it with $\chi^2_{20,0.95}/20\approx1.57$: find the smallest detectable $\delta$ at $v=15$~m/s and at $v=2$~m/s. (e) Interpret the result for urban versus motorway driving and propose how $\delta$ could be estimated rather than merely detected. The script \repopath{ch08\_\allowbreak p3\_\allowbreak timestamp\_\allowbreak offset.py} in \repopath{companion\_\allowbreak repo/\allowbreak exercises/\allowbreak ch08/\allowbreak solutions/\allowbreak } simulates the scenario.
\end{prob}

\begin{prob}\label{prob:ch08:4}
\textbf{Sim-to-real transfer protocol (design).} A camera--LiDAR detector with learned covariances has been trained entirely in simulation. Propose a protocol, using only unlabelled real recordings, to decide modality by modality whether the detector and its uncertainty model transfer. Specify the statistics (for instance the NIS of the tracker fed by the detector, the conflict mass between the camera and LiDAR branches, and the temporal consistency of detections), the reference values expected if transfer were perfect, and accept/reject thresholds. Discuss what a small labelled subset would add.
\end{prob}

\begin{prob}\label{prob:ch08:5}
\textbf{Cooperative perception message (discussion).} A roadside unit broadcasts object tracks to approaching vehicles. (a) List the minimal fields a message must carry for the receiver to apply covariance intersection \eqref{eq:ci-est}--\eqref{eq:ci-cov} and evidence discounting consistently. (b) Explain what goes wrong if the covariance is omitted and the receiver assumes a default one, and if the same track is received through two paths without provenance. (c) Propose how the receiver could estimate its trust in the sender from the consistency of past messages and how that trust should enter the fusion.
\end{prob}

\begin{prob}\label{prob:ch08:6}
\textbf{A foundation model as an information source (discussion).} A multimodal language model labels rare objects in camera images with a free-text class and a self-reported confidence. Explain how its output can be mapped to a basic belief assignment on a frame of discernment, how a discount factor can be estimated from a held-out set, how the conflict mass $K$ with the explicit sensors should be interpreted, and where in the pipeline of Fig.~\ref{fig:hybrid-pipeline} such a source should be inserted so that it cannot override a consistent classical estimate.
\end{prob}
\end{addedblock}

% ============================================
